% C44_POS, arithmetic continuation, /root/fermat_aj.
% GPT-6 Astra Pro; effort setting not exposed. 2026-11-22 (Europe/Berlin).
% New proofs below are subject to the recorded whole-proof check.

\section{The Fermat regulator: an integral projector and the exact
unramified obstruction}

The notation in this section is that of JAC1964 JAC-P13, JAC-P17,
AR-16a and AR-23.  In particular, $J/\mathbf Q$ is the Jacobian of
$x^4+y^4=z^4$, $N=h^3(J)^G$ with $G=\mu_4^2$, and
$\eta\in V:=\operatorname{CH}^2(N)_{\mathbf Q}$ is the specified
Fourier--Ceresa class.  The nonzero vector
$\alpha_5(\eta)$ spans
$S_5=H^1_f(\mathbf Q,H^3(N_{\overline{\mathbf Q}},\mathbf Q_5)(2))$
by the carried input AR-23, using AR-19 and H9-T15.  We do not reprove
that Selmer calculation, nor its nonzero cycle.  None of the deductions
in this section uses 032 or Milne's report.

The question $\ker\alpha_5=0$ remains unresolved.  The findings here
locate its integral obstruction more precisely.  They do not identify
a Selmer dimension with a Chow dimension.

\subsection{An integral model at the prime 5}

\begin{proposition}[C44-F01]\label{C44-F01}
Let $R=\mathbf Z_{(5)}$ and $S=\operatorname{Spec}\mathbf Z[1/10]$.
Let $\mathcal J/S$ be the Jacobian of the smooth Fermat curve
$\mathcal C:x^4+y^4=z^4$.  There is an exact idempotent, with
coefficients in $\mathbf Z[1/2]\subset R$, given by
\[
 e=\frac{1-[-1]^*}{2}\,\frac1{16}\sum_{g\in G}g^*.
 \tag{F01.1}
\]
The sum is formed after $S'=\operatorname{Spec}\mathbf Z[i,1/10]$
and descended to $S$.  Its rational generic-fibre motive is exactly
$N$.  Put
\[
 A=e\operatorname{CH}^2(J;R).
\]
Then restriction is an isomorphism
\[
 e\operatorname{CH}^2(\mathcal J;R)
       \xrightarrow{\ \sim\ } A,
 \qquad A\otimes_R\mathbf Q=V.                 \tag{F01.2}
\]
Moreover $e$ acts on geometric $5$-adic cohomology only in degree
three, where its image is a free rank-two direct summand.  Thus
\[
 T=eH^3_{\mathrm{et}}(J_{\overline{\mathbf Q}},\mathbf Z_5)(2)
 \tag{F01.3}
\]
is an actual $\mathbf Z_5$ lattice for the representation in the task,
and $eH^b_{\mathrm{et}}(J_{\overline{\mathbf Q}},\mathbf Z/5^n(2))=0$
for $b\ne3$ and every $n\ge1$.
\end{proposition}

\begin{proof}
The Fermat equation is smooth over $\mathbf Z[1/2]$: a singular
projective geometric point would have $x=y=z=0$, since all three
partial derivatives are units times the corresponding cubes.
The section $(1:0:1)$ gives its relative Jacobian, an abelian scheme.
The group action and inversion extend over $S$.  The two averaging
operators in (F01.1) commute and are idempotent.  They use only
powers of two as denominators.  Galois conjugation on $S'/S$ permutes
the sixteen group elements.  Restriction and norm, whose composite
is multiplication by two, descend the idempotent with the asserted
coefficients.  This descent is valid for Chow correspondences and
also for their actions on the cohomology theories used below.

The six characters of $G$ on $H^1(J)$ are
\[
 (1,1),(1,2),(2,1),(3,3),(3,2),(2,3)\quad\text{in }(\mathbf Z/4)^2;
 \tag{F01.4}
\]
these are the carried calculation in JAC-P17.  In their exterior
powers, an odd-degree invariant occurs only in degree three.  The
two invariant triples are the first three and the last three
characters.  Their types are $(3,0)$ and $(0,3)$.  The exact rational
identification with the projector of $N$ is JAC-P13/JAC-P17; in
particular (F01.1) is an integral-at-5 representative of the same
rational summand, not a new motive inferred solely from its
realization.  More explicitly, averaging is zero on $h^1(J)$
because it is induced by the zero norm endomorphism
$\sum_{g\in G}g$.  It is zero on $h^5(J)$ by the duality
identification through the $G$-invariant principal polarization.
Inversion removes the even Deninger--Murre pieces.  Thus the
rational idempotent on the full motive equals $\pi_3p_G$.
The geometric cohomology of an abelian variety is the
exterior algebra of its free degree-one cohomology.  Because the
projector has denominators prime to five, its integral image is a
direct summand and has no torsion.  The character calculation
therefore proves the assertions with both $\mathbf Z_5$ and
$\mathbf Z/5^n$ coefficients.

We next verify that $e$ annihilates the vertical terms in Chow
localization.  On any fibre of $\mathcal J$, the operator
$(1-[-1]^*)/2$ sends divisor classes into
$\operatorname{Pic}^0\otimes R$.  The endomorphism
$\sum_{g\in G}g$ of $\mathcal J_{S'}$ is zero.  Indeed its generic
fibre has zero tangent map by (F01.4), so is zero, and equality of
endomorphisms of an abelian scheme is detected on its dense generic
fibre.  Its dual is also zero, which says that
$\sum g^*$ annihilates $\operatorname{Pic}^0$.  Consequently
\[
 e\operatorname{CH}^1(\mathcal J_s;R)=0
 \quad\text{for every closed }s\in S.            \tag{F01.5}
\]
When a fibre does not contain $i$, check this after the extension
of degree at most two and descend using restriction--norm.

Chow localization gives
\[
 \bigoplus_{s\in S^{(1)}}\operatorname{CH}^1(\mathcal J_s;R)
 \longrightarrow\operatorname{CH}^2(\mathcal J;R)
 \longrightarrow\operatorname{CH}^2(J;R)\longrightarrow0.
 \tag{F01.6}
\]
The action of every automorphism in (F01.1) commutes with these
maps.  Applying the exact direct-summand functor $e$ and (F01.5)
proves (F01.2).  Notice that this argument is integral at five;
it does not invoke rational vanishing of the whole Chow group of
a special fibre.
\end{proof}

\paragraph{Locators and scope.}
The localization sequence is Fulton, \emph{Intersection theory},
Proposition~1.8 and its arithmetic-base version in \S20.2.
The Picard-scheme construction is FGA, Expos\'e~232,
Theorem~3.1 (printed pp.~148--149), applied to the smooth
projective relative curve; the canonical abelian component is
FGA, Expos\'e~236, Theorem~3.5 (printed p.~239), with the smooth
Picard schemes of the curve fibres.  The latter fibre smoothness
also follows from the vanishing of their $H^2(\mathcal O)$.
The relative-motive formalism over schemes smooth over a Dedekind
domain is recalled in the supplied Laga--Shnidman source,
\emph{Ceresa cycles of bielliptic Picard curves}, \S4.1--4.5;
the extension of Deninger--Murre to that base is explained at the
start of \S4.2, citing K\"unnemann, \S1, Remark~1.1.
The generic projector, the character list, and its rational field
of definition are JAC1964 JAC-P13 and JAC-P17.  The new assertion
here is the integral-at-5 localization isomorphism for that
projector.  The finite character check is reproduced by
\texttt{work/arithmetic/check\_fermat\_projector.py}; the script
does not prove any statement about rational equivalence.

\begin{lemma}[C44-F02]\label{C44-F02}
For the lattice $T$ in (F01.3),
\[
 H^0(\mathbf Q,T/5^nT)=0\quad(n\ge1),
 \qquad H^0(\mathbf Q,T\otimes\mathbf Q_5/\mathbf Z_5)=0.
 \tag{F02.1}
\]
Consequently
\[
 eH^3_{\mathrm{et}}(J,\mathbf Z/5^n(2))=0,
 \qquad
 eH^4_{\mathrm{et}}(J,\mathbf Z/5^n(2))
       \simeq H^1(\mathbf Q,T/5^nT).             \tag{F02.2}
\]
\end{lemma}

\begin{proof}
We make AR-23's Frobenius-at-11 check independent of a choice of
integral lattice.  By JAC-P17, the weight-three character is
$\psi_+^3$, where $\psi_+$ is the character of
$E_+:y^2=x^3+x$.  Direct counting gives
$\#E_+(\mathbf F_{11})=12$, so the degree-one Frobenius
polynomial is $X^2+11$.  Cubing its two roots gives
$X^2+11^3$ on $N$.  After the twist $(2)$, the polynomial is
\[
                 X^2+\frac1{11}.
\]
It follows that
\[
       \det(1-\operatorname{Fr}_{11}\mid T)
              =\frac{12}{11}\in\mathbf Z_5^\times.
 \tag{F02.3}
\]
For any stable $\mathbf Z_5$ lattice, this determinant condition
makes $1-\operatorname{Fr}_{11}$ invertible.  Its reductions are
invertible modulo every $5^n$, and it is invertible on the torsion
quotient as well.  Invariance under the entire Galois group implies
invariance under this Frobenius, proving (F02.1).  The conclusion
does not change if arithmetic rather than geometric Frobenius is
used, since both Frobenius and its inverse are integral units.

Apply the Hochschild--Serre spectral sequence with finite
coefficients and the integral projector $e$.  By C44-F01 the
projected geometric cohomology occurs only in row $b=3$.
The spectral sequence therefore degenerates on that row and
identifies its total degrees three and four with the two Galois
groups in (F02.1)--(F02.2).
\end{proof}

\paragraph{Check scope.}
The numerical part of (F02.3), including the point count and its
reduction $12/11\equiv2\pmod5$, is an exact integer/rational
computation in the accompanying script.  This is a lattice-uniform
form of the Frobenius test used in AR-23, not a new proof of
AR-23's nonzero Abel--Jacobi class.

\subsection{The precise obstruction after passing to finite coefficients}

Write
\[
 B=eH^4_{\mathrm{et}}(J,\mathbf Z_{(5)}(2)),
 \qquad
 C=eH^3_{\mathrm{nr}}(J,\mathbf Q_5/\mathbf Z_5(2)).
 \tag{F03.1}
\]
The cohomology defining $B$ is \emph{\'etale motivic cohomology},
not continuous $5$-adic cohomology.  The distinction is essential.
The group $C$ is unramified cohomology over the ground field
$\mathbf Q$, with residues at the codimension-one points of $J$.
The operations defining $e$ act on it by restriction, norm and
the automorphisms of $J_{\mathbf Q(i)}$.

\begin{theorem}[C44-F03]\label{C44-F03}
There is a canonical short exact sequence
\[
                 0\longrightarrow A\longrightarrow B
                   \longrightarrow C\longrightarrow0,
 \tag{F03.2}
\]
and both $A$ and $B$ have no nonzero $5$-primary torsion.
For every $n\ge1$, the kernel of the finite-coefficient cycle map
\[
 c_n:A/5^nA\longrightarrow H^1(\mathbf Q,T/5^nT)
\]
is canonically $C[5^n]$.  In particular, for
$\widehat A=\varprojlim_n A/5^nA$, the completed cycle map satisfies
\[
 \ker\!\left(\widehat c:\widehat A\longrightarrow
                    H^1_{\mathrm{cont}}(\mathbf Q,T)\right)
          = T_5C:=\varprojlim_n C[5^n],            \tag{F03.3}
\]
where the transition on $C[5^n]$ is multiplication by five.

Let
\[
 \rho:A\otimes_R\mathbf Q_5\longrightarrow
                  \widehat A\otimes_{\mathbf Z_5}\mathbf Q_5,
 \qquad K_{\mathrm{comp}}=\ker\rho.
\]
The full scalar-extended Abel--Jacobi kernel in the question has
the exact description
\[
 0\longrightarrow K_{\mathrm{comp}}
 \longrightarrow\ker\alpha_5
 \xrightarrow{\rho}
 \operatorname{im}\rho\cap
             (T_5C\otimes_{\mathbf Z_5}\mathbf Q_5)
 \longrightarrow0.                              \tag{F03.4}
\]
The intersection is taken inside
$\widehat A\otimes\mathbf Q_5$.  Neither group on the two ends
of (F03.4) has been proved to vanish.
\end{theorem}

\begin{proof}
For a smooth variety over a field, the codimension-two
comparison sequence is
\[
 0\longrightarrow\operatorname{CH}^2(J)
 \longrightarrow H^4_{\mathrm{et}}(J,\mathbf Z(2))
 \longrightarrow H^3_{\mathrm{nr}}(J,\mathbf Q/\mathbf Z(2))
 \longrightarrow0.
\]
The exact locator is Kahn, \emph{On the injectivity and
non-injectivity of the $l$-adic cycle class maps}, (1.5),
author p.~3, citing his \emph{Classes de cycles motiviques
\'etales}, Proposition~2.9.  Localizing at five and applying
the idempotent $e$ gives (F03.2).

The coefficient triangle for $\mathbf Z_{(5)}(2)$ has
$\mathbf Z/5^n(2)=\mu_{5^n}^{\otimes2}$ as its finite-coefficient
term.  Its long exact sequence and (F02.2) give
$B[5^n]=0$ and an injection
\[
 B/5^nB\hookrightarrow
       eH^4_{\mathrm{et}}(J,\mathbf Z/5^n(2)).     \tag{F03.5}
\]
The relevant coefficient comparison is the Merkurjev--Suslin
case of the norm-residue theorem; it is recorded in Kahn,
loc.~cit., (3.2) and (4.1), author pp.~7 and~10.
Since $A\hookrightarrow B$, also $A[5^n]=0$.

Apply multiplication by $5^n$ to (F03.2).  The snake lemma gives
\[
 0\longrightarrow C[5^n]\xrightarrow{\delta_n}A/5^nA
 \longrightarrow B/5^nB.
 \tag{F03.6}
\]
Explicitly, for $c\in C[5^n]$, choose $b\in B$ above $c$;
then $5^n b\in A$, and
$\delta_n(c)=5^n b\bmod5^nA$.  This is independent of the
choice of lift.  The absence of $5^n$-torsion in $B$ proves
its injectivity.  Composing with (F03.5) and the
Hochschild--Serre identification gives precisely $c_n$.
Thus (F03.6) proves the finite-coefficient assertion.

If $c\in C[5^{n+1}]$ and $b$ is a lift, then the lift $5b$
of $5c$ shows that reduction of $\delta_{n+1}(c)$ modulo $5^n$
equals $\delta_n(5c)$.  The transition is therefore multiplication
by five.  Taking inverse limits in (F03.6) is exact at the first
two terms, so the kernel of the completed finite-coefficient
cycle map is $T_5C$.  This uses left exactness only; no
unverified surjectivity of transitions is needed.

Continuous cohomology has the comparison exact sequence
\[
 0\longrightarrow\varprojlim{}^1_n
       H^0(\mathbf Q,T/5^nT)
 \longrightarrow H^1_{\mathrm{cont}}(\mathbf Q,T)
 \longrightarrow\varprojlim_n H^1(\mathbf Q,T/5^nT)
 \longrightarrow0.
\]
Its left term is zero by C44-F02.  Equivalently, in total
degree four on $J$ the potential $\varprojlim{}^1$ term is the
limit of the zero groups in (F02.2).  Thus the completed map
just constructed is the continuous cycle map, proving (F03.3).
This is why an inverse limit is legitimate in this particular
summand although it is not interchangeable with continuous
cohomology for an arbitrary variety over a number field.

After tensoring by $\mathbf Q_5$, the cycle map from $A$
factors as $\widehat c\circ\rho$.  Its rational projected
Hochschild--Serre component is the Abel--Jacobi map
$\alpha_5$ from the task.  The inclusion of its image in
$H^1_f$ is the carried AR-16a formalism and does not change
the kernel.  Flatness of $\mathbf Q_5$ over $\mathbf Z_5$
identifies the kernel of $\widehat c\otimes\mathbf Q_5$ with
$T_5C\otimes\mathbf Q_5$.  Taking the inverse image under
$\rho$ and then its image yields (F03.4).
\end{proof}

\paragraph{What this attack removes, and what it leaves.}
C44-F01 removes the ambiguity of a varying arithmetic spread for
the Chow source.  C44-F02 removes both the global degree-three
term and the continuous-cohomology $\varprojlim{}^1$ ambiguity
from the finite-coefficient calculation.  The remaining group
in (F03.3) is the actual Tate module of one specified unramified
group.  It is not the Selmer group.  A Selmer bound does not bound
this Tate module or prove the tensor-to-completion injectivity in
(F03.4).  Specialization of rational odd-motive Chow classes,
already ruled out as a detector by AR-20, supplies neither
missing assertion.

\subsection{A bounded-denominator formulation on rational planes}

For a finite free $R$-submodule $M\subset A$, let
$W_M=M\otimes_R\mathbf Q\subset V$.  Equation (F03.2) and
the absence of $5$-torsion identify $B$ with a subgroup of $V$:
if $b\in B$, choose $5^n b=a\in A$ and send $b$ to $a/5^n$.
This map is well defined and injective, since $C$ is $5$-primary
torsion.  This convention makes $B\cap W_M$ unambiguous.

\begin{proposition}[C44-F04]\label{C44-F04}
The following are equivalent for every finite free $M\subset A$:
\begin{enumerate}
\item $\alpha_5$ is injective on $M\otimes_R\mathbf Q_5$;
\item $T_5(B/M)=0$;
\item $(B\cap W_M)/M$ is finite;
\item there is an integer $b_M\ge0$ such that
      $5^{b_M}(B\cap W_M)\subset M$.
\end{enumerate}
More precisely,
\[
 \ker(\alpha_5|_{M\otimes\mathbf Q_5})
       \simeq T_5(B/M)\otimes_{\mathbf Z_5}\mathbf Q_5.
 \tag{F04.1}
\]
Choose a nonzero multiple $a_\eta\in A$ of the carried class
$\eta$.  It is enough to verify these equivalent conditions for
the modules
\[
                  M_z=R a_\eta+Rz\quad(z\in A).
 \tag{F04.2}
\]
The condition holds when $z$ lies in $\mathbf Q\eta$, by AR-23.
For arbitrary $z$ the bounded-denominator condition in (4) is
unproved.
\end{proposition}

\begin{proof}
Apply the argument of (F03.6) to
$0\to M\to B\to B/M\to0$.  Both $M$ and $B$ have no
$5$-torsion.  It gives
\[
 0\longrightarrow T_5(B/M)\longrightarrow\widehat M
                         \longrightarrow\widehat B.
\]
Since $M$ is finite free over $R$, its completion is
$M\otimes_R\mathbf Z_5$.  The injection (F03.5) implies
that the kernel of the map to $\widehat B$ is exactly the
kernel of the continuous cycle map.  Rationalizing proves
(F04.1).

All $5$-power-torsion elements of $B/M$ lie in
\[
 D_M=(B\cap W_M)/M\ \subset\ W_M/M
          \simeq(\mathbf Q_5/\mathbf Z_5)^{\operatorname{rank}M}.
\]
Conversely every element of $D_M$ is $5$-power torsion, since
$R$ already inverts the other primes.  Thus $T_5(B/M)=T_5D_M$.
A subgroup of a finite direct sum of $\mathbf Q_5/\mathbf Z_5$
is a direct sum of finitely many copies of
$\mathbf Q_5/\mathbf Z_5$ and a finite group.  One proof is to
take its Pontryagin dual: it is a quotient of the finite free
$\mathbf Z_5$ module, hence is finitely generated; the structure
theorem over the discrete valuation ring $\mathbf Z_5$ and
duality give the assertion.  It follows that $T_5D_M$ is finite
free, and vanishes exactly when $D_M$ is finite.  Finiteness of
this subgroup is equivalent to a bound on its exponent, giving
the equivalence with (4) and also with (1), by (F04.1).

Finally AR-23 makes $\alpha_5(a_\eta)\ne0$ and gives a
one-dimensional target.  If all the two-generator modules
(F04.2) have injective scalar-extended map, none can have rank
two.  Hence every $z\in A$ belongs to $\mathbf Q\eta$, so
$V=\mathbf Q\eta$ and the full map is injective.  The converse
is immediate.  On the line $\mathbf Q\eta$, AR-23 already gives
injectivity, so the equivalent bounded-denominator conclusion
is proved on that line without any hypothesis about the rest
of the Chow group.
\end{proof}

\begin{proposition}[C44-F05: an algebraic model of the missing step]
\label{C44-F05}
Even finite generation of the integral Chow-like source,
one-dimensional surjective regulator image, a nonzero distinguished
class, and a positive definite real form do not formally imply
the scalar-extended injectivity in C44-F04.  There is an explicit
model satisfying all those properties in which the completed
kernel is one-dimensional and arises from a nonzero unramified-like
Tate module.
\end{proposition}

\begin{proof}
Choose the root $a\in\mathbf Z_5$ of $a^2=6$ with $a\equiv1\pmod5$.
Hensel's lemma gives this root, and $a\notin\mathbf Q$ because
six is not a rational square.  Put
\[
 A_0=R^2,\qquad
 B_0=\{(r,s)\in\mathbf Q^2:r+as\in\mathbf Z_5\},
 \qquad f(r,s)=r+as.
\]
Then $A_0\subset B_0$, both are $5$-torsion-free, and
\[
  B_0/A_0\simeq\mathbf Q_5/\mathbf Z_5.
 \tag{F05.1}
\]
To see the last statement, take the kernel of the surjection
$(\mathbf Q_5/\mathbf Z_5)^2\to\mathbf Q_5/\mathbf Z_5$,
$(r,s)\mapsto r+as$, whose restriction on the first coordinate
is the identity.  The inverse image of this kernel in
$\mathbf Q^2$ is exactly $B_0$.

The scalar-extended regulator is
$\mathbf Q_5^2\to\mathbf Q_5$, $(r,s)\mapsto r+as$.
It is surjective, sends $\eta_0=(1,0)$ to one, and has kernel
$\mathbf Q_5(-a,1)$.  It is even injective on the underlying
rational vector space $\mathbf Q^2$.  The form
$h((r,s),(r',s'))=rr'+ss'$ is positive definite and
$h(\eta_0,\eta_0)=1$.  Neither fact removes the scalar-extended
kernel.

For a completely explicit division tower, choose integers $a_n$
with $a_n\equiv a\pmod{5^n}$ and
$a_{n+1}\equiv a_n\pmod{5^n}$, and set
\[
 b_n=5^{-n}(-a_n,1)\in B_0.
\]
Then $b_n\bmod A_0$ has order $5^n$ and
$5b_{n+1}-b_n\in A_0$.  The resulting nonzero element of
$T_5(B_0/A_0)$ maps to $(-a,1)\in\widehat A_0$.
Thus the precise obstruction in C44-F03 is present in the
model, despite all the other listed data.  This is a model of
an invalid inference, not a geometric counterexample to
Abel--Jacobi injectivity or to a height conjecture.
\end{proof}

\subsection{What finite coefficients see of the fourfold component}

\begin{proposition}[C44-F06]\label{C44-F06}
Let $F/\mathbf Q$ be finite and let $D/F$ be a smooth projective
geometrically connected curve with a point $d_0\in D(F)$.
On $X=D\times J_F$, put
\[
 q=\pi_1^D\otimes e,\qquad
 A_D=q\operatorname{CH}^2(X;R),\qquad
 B_D=qH^4_{\mathrm{et}}(X,R(2)),
 \quad R=\mathbf Z_{(5)}.
\]
Then $B_D$ is uniquely divisible by five and hence is a
$\mathbf Q$-vector space.  The motivic comparison identifies it
canonically with the rational Chow component:
\[
       B_D\simeq A_D\otimes_R\mathbf Q.
                                                        \tag{F06.1}
\]
In particular $A_D$ has no $5$-primary torsion, and the
finite-coefficient cycle map on $A_D/5^nA_D$ is zero for every $n$.
Writing
\[
 C_D=qH^3_{\mathrm{nr}}(X,\mathbf Q_5/\mathbf Z_5(2)),
\]
one has
\[
 C_D=B_D/A_D,\qquad
 C_D[5^n]\simeq A_D/5^nA_D,
 \qquad T_5C_D\simeq\widehat A_D.             \tag{F06.2}
\]
Thus finite generation of $B_D$ as an $R$-module would force
the entire rational component in JAC-P18 to vanish.  Such finite
generation is not established here; finite generation of $A_D$
alone would not have that consequence.
\end{proposition}

\begin{proof}
The point projector $\pi_1^D$ is integral.  By C44-F01 and
the K\"unneth formula, $q$ selects geometric cohomology only in
degree four, where the integral lattice after twist is
\[
 T_D=H^1_{\mathrm{et}}(D_{\overline F},\mathbf Z_5)
                  \otimes_{\mathbf Z_5}T.
\]
It is free, of rank $4g(D)$.  At any place of $F$ over five,
its rational realization has Hodge--Tate weights
$-2,-1,1,2$ (each with multiplicity $g(D)$, up to the
simultaneous opposite weight convention).  Indeed the two
factors before twist have weights $\{0,1\}$ and $\{0,3\}$.
None is zero.  The Hodge--Tate comparison for smooth proper
varieties therefore implies
\[
                 H^0(F,T_D)=0.                 \tag{F06.3}
\]
Here one may use Scholze, \emph{$p$-adic Hodge theory for
rigid-analytic varieties}, Corollary~1.8, author p.~4, applied
to the analytification, together with functoriality for the
algebraic projectors.  This use needs neither good reduction
of $D$ at five nor a Hodge or Tate conjecture.

Finite-coefficient Hochschild--Serre now gives
\[
 qH^3_{\mathrm{et}}(X,\mathbf Z/5^n(2))=0,
 \qquad
 qH^4_{\mathrm{et}}(X,\mathbf Z/5^n(2))
       =H^0(F,T_D/5^nT_D).
                                                        \tag{F06.4}
\]
As in C44-F03, the coefficient triangle implies that
$B_D[5^n]=0$ and injects $B_D/5^nB_D$ into the second group
in (F06.4).  Taking inverse limits gives
\[
 \widehat B_D\hookrightarrow
        \varprojlim_n H^0(F,T_D/5^nT_D)
        =H^0(F,T_D)=0.
\]
The maps $B_D/5^{n+1}B_D\to B_D/5^nB_D$ are surjective,
so each projection from their inverse limit is surjective.
Consequently $B_D/5^nB_D=0$ for every $n$.  Together with
$B_D[5^n]=0$ this proves unique divisibility by five.
The other primes are already inverted in $R$.

The comparison sequence (F03.2) is valid in codimension two
on $X$ as well, and is compatible with $q$.  Its quotient
is $5$-primary torsion.  Rationalizing its injection
$A_D\hookrightarrow B_D$, and using that $B_D$ is already
a rational vector space, proves (F06.1).  The finite-coefficient
cycle map factors through $B_D/5^nB_D=0$.  Applying the snake
lemma to multiplication by $5^n$ gives (F06.2), including the
transition by multiplication by five; its inverse limit is
the final displayed isomorphism.

A nonzero rational vector space is not finitely generated as
a $\mathbf Z_{(5)}$-module.  Thus the stated finite generation
condition forces $B_D=0$.  JAC-P18 has already proved that
every algebraic class in this projected component is
homologically zero, so this is the vanishing that its
componentwise Murre~(D) would require.  The elementary inclusion
$R\subset\mathbf Q$ shows why finite generation of the smaller
group $A_D$ cannot substitute for finite generation of $B_D$.
\end{proof}

\paragraph{Scope of the additional attack.}
C44-F06 does not prove Murre~(D).  It identifies the relevant
\'etale motivic group as entirely divisible, with all of its
completed Chow source accounted for by unramified cohomology.
This makes precise why torsion-freeness or finite-coefficient
cycle computations alone cannot kill the JAC-P18 obstruction.
It concerns cote~44, Th.~1--3 (scan~002) and the boundary of
Th.~5 (scan~004); it does not change the known polarization
on the incidence higher Picard factor.  In the frameworks of
SGA and Chow motives, it is the tensor summand
$h^1(D)\otimes N$ of weight four with no $(2,2)$ part,
followed by the same codimension-two \`etale-motivic comparison
as C44-F03.  Its vanishing is a rational-equivalence issue,
not the semisimplicity or signed $I$ assertion of cote~16,
scan~049.  No 032/Milne input is used.

\subsection{The remaining finiteness question on one arithmetic model}

\begin{proposition}[C44-F07]\label{C44-F07}
For the abelian scheme $\mathcal J/S$ of C44-F01, restriction
also gives an isomorphism on the projected \emph{\'etale
motivic} groups:
\[
 eH^4_{\mathrm{et}}(\mathcal J,R(2))
                 \xrightarrow{\ \sim\ } B.
                                                        \tag{F07.1}
\]
The completion $\widehat B$ is a free $\mathbf Z_5$-module
of rank exactly one.  Its rational realization is the Selmer
line already realized by $\eta$.

The following two conditions are equivalent:
\begin{enumerate}
\item the scalar-extended Fermat Abel--Jacobi map is injective;
\item $eH^4_{\mathrm{et}}(\mathcal J,R(2))$ is finitely
      generated over $R$.
\end{enumerate}
Condition (2) is not proved here.  In particular the proved
finite generation of the completion is not being substituted
for it.
\end{proposition}

\begin{proof}
Purity and localization for \'etale motivic cohomology
give the exact segment
\[
 \bigoplus_{s\in S^{(1)}}
 H^2_{\mathrm{et}}(\mathcal J_s,R(1))
 \longrightarrow H^4_{\mathrm{et}}(\mathcal J,R(2))
 \longrightarrow H^4_{\mathrm{et}}(J,R(2))
 \longrightarrow\bigoplus_{s\in S^{(1)}}
 H^3_{\mathrm{et}}(\mathcal J_s,R(1)).
                                                        \tag{F07.2}
\]
An exact source for the needed purity is Geisser,
\emph{Motivic cohomology over Dedekind rings},
Theorem~1.2(1), published p.~774: after inverting the residue
characteristic no truncation is needed.  All residue
characteristics of $S$ differ from five and are units in $R$.
The Bloch--Kato premise of that 2004 statement is supplied by
the norm-residue theorem (in this weight, Merkurjev--Suslin).
The localization of the cycle complex and passage to the
generic fibre give the displayed direct sum.  This is also
the localization used in Geisser, \emph{Hasse principles for
\'etale motivic cohomology}, \S4, diagram~(2).

The first outer group in (F07.2) is a sum of Picard groups,
because $\mathbf Z(1)_{\mathrm{et}}=\mathbf G_m[-1]$.
Its $e$-part vanishes by C44-F01.  The last is a sum of
cohomological Brauer groups localized at five.  For a fibre
$\mathcal J_s$, finite-coefficient Hochschild--Serre and the
character calculation of C44-F01 give
\[
           eH^2_{\mathrm{et}}(\mathcal J_s,\mu_{5^n})=0
           \quad(n\ge1).
\]
Indeed its projected geometric cohomology starts in degree
three.  Kummer's sequence implies
$e\operatorname{Br}(\mathcal J_s)[5^n]=0$ for every $n$.
The Brauer group is torsion, so its localization at five has
zero $e$-part.  Apply $e$ to (F07.2) to obtain (F07.1).

The lattice $T$ is a lisse $\mathbf Z_5$-sheaf on $S$.
Proper smooth base change and the same single-row calculation
give
\[
 eH^4_{\mathrm{et}}(\mathcal J,\mathbf Z/5^n(2))
          =H^1_{\mathrm{et}}(S,T/5^nT),
 \qquad eH^3_{\mathrm{et}}(\mathcal J,\mathbf Z/5^n(2))=0.
                                                        \tag{F07.3}
\]
For the second identity, Frobenius at eleven still belongs
to $\pi_1(S)$, so C44-F02 applies.  Finite-coefficient
cohomology of this arithmetic scheme is finite.  Alternatively,
in degree one this follows from the finiteness of number
field extensions of bounded degree unramified outside
$\{2,5,\infty\}$.
It follows from the coefficient comparison that each
$B/5^nB$ is finite, and that
\[
 \widehat B\hookrightarrow
           H^1_{\mathrm{cont}}(S,T).
                                                        \tag{F07.4}
\]
There is no inverse-limit correction by (F07.3).
The group on the right is a finite free $\mathbf Z_5$-module:
it is a compact inverse limit of finite groups, its quotient
modulo five injects into $H^1(S,T/5T)$, and hence compact
Nakayama gives finite generation; its $5$-torsion vanishes
because $H^0(S,T/5T)=0$.  Restriction to
$H^1_{\mathrm{cont}}(\mathbf Q,T)$ is injective by inflation.

The group $B$ is rationally $V$, by (F03.2).  Therefore the
rational image of every element of $B$ in (F07.4) belongs to
the Selmer line of AR-23.  Its intersection with the finite
free lattice $H^1_{\mathrm{cont}}(S,T)$ is a closed rank-one
lattice.  The image of $B$ is dense in $\widehat B$, whose
image in (F07.4) is closed by compactness.  Hence
$\widehat B$ lies in that rank-one lattice.  It is nonzero,
because $a_\eta\in A\subset B$ has nonzero regulator.
This proves its asserted rank, including its rational image.

If $B$ is finitely generated over $R$, torsion-freeness makes
it finite free.  Its rank equals the rank of $\widehat B$,
namely one.  As $A\subset B$ and $A\otimes_R\mathbf Q=V$,
this forces $V=\mathbf Q\eta$, and AR-23 then proves
injectivity.  Conversely, injectivity and AR-23 give
$V=\mathbf Q\eta$.  Apply C44-F04 to $M=Ra_\eta$: since
its rational span is all of $V$, the bounded-denominator
conclusion places $B$ inside $5^{-b_M}M$.  Thus $B$ is a
submodule of a finite free $R$-module and is finitely
generated.  This proves the equivalence.
\end{proof}

\paragraph{Meaning of the arithmetic finiteness gate.}
C44-F07 makes the finiteness question a statement about one
specific summand on one scheme over $\mathbf Z[1/10]$.
Neither finite generation of ordinary Chow groups nor
finite generation of continuous \(5\)-adic cohomology proves
it.  Kahn, \S4, author pp.~10--11, explicitly isolates this
distinction; his Proposition~4.1 gives counterexamples to
unqualified \`etale-motivic finiteness on arithmetic schemes.
Geisser, \emph{Hasse principles}, Conjecture~4.3, states
Lichtenbaum's finiteness expectation for regular schemes
proper over $\mathbf Z$; that conjecture is not a theorem
for the present summand, and $\mathcal J/S$ is proper over
the indicated open base, not over all of $\mathbf Z$.
No claim of direct applicability or proof of that conjecture
is made.  In cote~44 this refines the arithmetic gate of
Th.~5, scan~004; in cote~16 it specifies information about
rational equivalence that conditions $C$ and $I$ on pure
realizations do not provide.  The equivalence does not by
itself count as a strictly smaller main gate under the
progress rule of the prompt.  No 032/Milne input is used.

\subsection{The three proposed arithmetic routes: exact scope}

\paragraph{C44-FA01: Euler systems.}
The supplied Lilienfeldt--Shnidman paper,
\emph{Derivatives of Rankin--Selberg $L$-functions and heights of
generalized Heegner cycles}, arXiv:2408.04375,
Theorem~1.1/13.4, identifies a height on a Hecke-generated cycle
space modulo the real-height radical.  Its Corollary~1.4 expressly
assumes the injectivity of its scalar-extended Abel--Jacobi map.
In \S15, author pp.~71--73, equation~(99) on p.~72 is the Euler-system
conclusion that a nonzero generalized Heegner cohomology class
spans the Selmer space; the next paragraph invokes the assumed
injectivity to pass back to Chow.  The displayed implication chain
on p.~73 marks both uses of Abel--Jacobi injectivity.  That is the exact missing
step, and C44-F03--F04 describe its integral obstruction for the
Fermat summand.  No Euler-system relation at the level of
Galois cohomology bounds the denominators in (F04.2).

There are also independent obstacles to a direct substitution
in the displayed theorem: its standing assumptions in \S1.1
require an imaginary quadratic field of odd discriminant and
an unramified Hecke character, whereas the Fermat character
over $\mathbf Q(i)$ has conductor $(1+i)^4$ by JAC-P17/AR-17.
The stated range is $0<t\le k-1$, so the weight-four central
case with $k+t=2$ is not a direct instance.  These scope failures
are not the reason the existing Selmer result fails---that result
is already available as AR-19/H9-T15---but they prevent the paper
from being treated as a new Chow-generation theorem here.

The Loeffler--Zerbes constructions consulted in this run
interpolate or bound Galois-cohomology classes, not the entire
Chow group.  The exact example is Jetchev--Loeffler--Zerbes,
\emph{Heegner points in Coleman families}, Theorem~A
(published p.~125), equivalently Theorem~5.4.1, whose
specializations are \emph{Abel--Jacobi images} of generalized
Heegner cycles.  An interpolation statement about these images
does not supply the missing denominator bound.  The existing
JAC1964 audit AR-17 and its
\texttt{LATEST\_RANK\_SOURCE\_SCOPE.tex} remain in force; their
statements are not reproved here.

\paragraph{C44-FA02: Gross--Schoen heights.}
AR-24--AR-25 are carried as the exact projected-diagonal and
sixteen-intersection identities.  Zhang,
\emph{Gross--Schoen cycles and dualising sheaves},
Theorem~2.3.5, author p.~27, proves the triple-intersection
identity; Conjecture~2.4.1, author p.~28, asserts semipositivity
of the tensor-correspondence pairing.  The latter is not proved
by the former.  Even adding that conjecture would give the
semipositivity conclusion recorded in AR-25, without an
equality criterion or the injectivity in task~3.  In particular
one must not substitute a real-height calculation for
the absent faithfulness of the $5$-adic regulator.  C44-F05
shows that positivity and the stated Selmer data alone are not
a formal bridge.  The supplied Gu draft, Theorem~1.3,
Corollary~5.4 and Conjecture~1.1, gives formulas and
conditional positivity, not that missing faithfulness.  Its
statements are treated here as supplied draft inputs, not as
independently verified public results, and none is used in the
proofs of C44-F01--F05.

\paragraph{C44-FA03: the algebraically trivial radical.}
JAC-P18 isolates the possible cycle
$Z=(\pi_1^C\otimes\pi_N)_*Z$ on $C\times J$, fixed by the
degree-four product projector.  Xu--Xu,
\emph{Remarks on Murre's conjecture on Chow groups},
Corollary~3.9, author p.~9, gives (A) and (B) for the product;
the earlier definition of (D), author p.~2, is the extra
cycle-class injectivity needed to annihilate this component.
Theorem (B) leaves degree four acting on codimension two and
therefore does not remove this component.  The special
Laga--Shnidman isomorphism in their Theorem~5.2 concerns a
different invariant summand of types $(2,1),(1,2)$, identified
with $h^1(E)(-1)$.  Such a Picard-motive isomorphism cannot hold
for the Fermat $N$, whose types are $(3,0),(0,3)$.
No proof of the needed componentwise (D) was found in these
sources or in the other primary sources checked in this run.

\paragraph{Programme crosswalk and current gate.}
C44-F01--F04 concern the arithmetic extension of cote~44,
scan~004, Th.~5, and the incidence/higher-Picard boundary of
Th.~1--3 on scan~002, as made explicit in the letter on
scans~008--009.  They do not alter JAC-T4.3 or establish a new
polarization of $P^i(X)$.  The actual cycle lives in the
extreme-type motive outside that Picard component.
In cote~16, scans~049--050, the conclusion is a limitation of
the passage from $C$, $I$, and semisimplicity of cohomological
correspondences to rational Chow groups: the kernel measured
by (F03.4) is invisible to those pure-realization assertions.
The corrected $A$, $A'$, $A''$ conditions of SN-C04 are
unchanged.

In Grothendieck's later language, C44-F01 is a statement over
one scheme over $\mathbf Z$; (F03.2) is the comparison from
codimension-two cycles to \emph{\'etale motivic} cohomology;
(F02.2)--(F03.4) concern $\ell$-adic sheaves, weights, and their
arithmetic realization.  For the generic fibre, SGA~6,
Expos\'e~XIV, \S4.1--4.2 (especially (4.1) and (4.4)),
identifies the rational Chow group in codimension two with
degree two of the topological and $\lambda$ filtrations on
$K_0(J)$.  The operators in (F01.1) act there by pullback.
This $\lambda$-ring translation supplies no bound on rational
equivalences.  No new EGA/FGA theorem or standard-conjecture
implication is asserted: the modern comparison theorems are
cited under their actual sources above.

The open rows \texttt{QFermat-AJ} and \texttt{Q35-full} are
\emph{not closed}.  Their exact new test is (F03.4), or,
equivalently, bounded five-power denominators in $B$ on every
rational plane (F04.2).  The original analytic order-one result
AR-11, positive value AR-13, and Selmer-line result AR-23
remain intact.  None of the present findings proves
$\operatorname{rank}\operatorname{CH}^2(N)=1$.
